\documentclass[11pt]{article}

\usepackage[margin=1.05in]{geometry}
\usepackage{amsmath,amssymb,amsthm}
\usepackage{booktabs}
\usepackage{enumitem}
\usepackage{xcolor}
\usepackage[colorlinks=true,linkcolor=blue!45!black,citecolor=blue!45!black,urlcolor=blue!45!black]{hyperref}

\newcommand{\R}{\mathbb{R}}
\newcommand{\code}[1]{\texttt{\small #1}}
\newcommand{\Hxx}{H_{\!X\!X}}
\newcommand{\xstar}{x^{\star}}

\theoremstyle{plain}
\newtheorem{proposition}{Proposition}
\theoremstyle{definition}
\newtheorem{condition}{Condition}
\theoremstyle{remark}
\newtheorem{remark}{Remark}

\title{\bfseries Conditional Gaussianity in the MAGI posterior:\\[2pt]
what an affine vector field does and does not buy}
\author{}
\date{\today}

\begin{document}
\maketitle

\begin{abstract}
\noindent
If the ODE's vector field is affine in the state at fixed parameters, MAGI's conditional
$p(X \mid \theta)$ is \emph{exactly} Gaussian. That claim invites an immediate objection: the
parameters may enter the field arbitrarily --- through exponentials, reciprocals, anything --- so
how can the posterior be Gaussian? The objection is sound and the claim survives, because the two
statements are about different objects. Gaussianity here is a property of the conditional in $X$;
$\theta$ enters only through coefficients that are constants with respect to $X$, so it can shape
the conditional's mean and covariance however it likes without changing the degree of the
polynomial in $X$. The joint and the $\theta$-marginal remain as non-Gaussian as the model makes
them.

We state and prove the proposition, exhibit a field that is affine in $X$ while depending on
$\theta$ through $e^{a}$, $\sin 3b$, $c^{2}$ and $(1+a^{2})^{-1}$, and verify numerically that its
log-density is exactly quadratic in $X$ at a displacement of $100$ state units while being
thoroughly non-quadratic in $\theta$. We then separate what the condition buys --- the elimination
of the states is exact, so a $d$-dimensional problem becomes a $p$-dimensional one with no
approximation error --- from what it does not: the remaining $p$-dimensional integral is not
Gaussian and still has to be done. Finally we measure how far the four benchmark systems sit from
the condition, which required first repairing the diagnostic that measures it: its perturbation was
absolute rather than relative, and on a system whose states reach $10^{5}$ that reported
$7.5\times10^{-7}$ where a like-for-like perturbation gives $1.1\times10^{-3}$.
\end{abstract}

\section{The objection}

The claim under examination, which appears in \cite{companion1} and in this codebase's
\code{condition\_A}, is:

\begin{quote}
If the vector field is affine in the state at fixed parameters, then $p(X \mid \theta)$ is exactly
Gaussian and no approximation is required to integrate the states out.
\end{quote}

The objection is that $\theta$ may enter $f$ in any way at all. Nothing in ``affine in the state''
restricts the $\theta$-dependence, so one could have $\dot x = -e^{\theta} x$, whose posterior in
$\theta$ is certainly not Gaussian. How can a posterior be exactly Gaussian when a parameter
appears inside an exponential?

It cannot, and the claim does not say it is. The claim is about the \emph{conditional}
$p(X \mid \theta)$ at a fixed value of $\theta$. Section~\ref{sec:prop} shows that $\theta$'s
arbitrariness is confined to coefficients that do not involve $X$, so it cannot disturb the
quadratic structure in $X$; Section~\ref{sec:notbuy} shows that it survives untouched in the
$\theta$-marginal, which is where the objection's intuition correctly applies.

\section{The log-density}
\label{sec:setup}

Write $n$ for the discretisation points, $D$ for the state components, $p$ for the parameters, and
$x = \operatorname{vec}(X) \in \R^{nD}$. MAGI's negative log-density at fixed observation noise is
a sum of four terms \cite{yang2021}:
\begin{equation}
\label{eq:U}
  U(\theta, x) = \tfrac12\Big[\,
    b\,(x-\mu)^\top G (x-\mu)
    \;+\; (x-y)^\top W (x-y)
    \;+\; b\, r^\top \Lambda\, r \,\Big]
  \;+\; \tfrac12 (\theta-\bar\theta)^\top P (\theta-\bar\theta),
\end{equation}
dropping terms independent of $(\theta, x)$. Here $b = \beta^{-1}$ is the prior tempering,
$G = \operatorname{blkdiag}_d(C_d^{-1})$ and $\Lambda = \operatorname{blkdiag}_d(K_d^{-1})$ are the
GP precision blocks, $W$ is diagonal with $\sigma_d^{-2}$ on observed entries and zero elsewhere,
$P$ is the parameter prior precision, and
\begin{equation}
\label{eq:resid}
  r(\theta, x) \;=\; F(x, \theta) - \dot\mu - M(x - \mu), \qquad
  M = \operatorname{blkdiag}_d(m_d),
\end{equation}
with $F(x,\theta)$ stacking $f(X_i, \theta, t_i)$ over grid points. Three observations about
\eqref{eq:U} carry the whole argument:

\begin{enumerate}[leftmargin=1.6em,itemsep=2pt]
\item The first two bracketed terms are quadratic in $x$ and involve $\theta$ \emph{not at all}.
\item The fourth term involves $\theta$ and involves $x$ not at all.
\item $\theta$ and $x$ meet in exactly one place: the residual $r$, and there only through $F$.
\end{enumerate}

So every question about the joint structure of \eqref{eq:U} is a question about $F$.

\section{The condition, and what it gives}
\label{sec:prop}

\begin{condition}[affine in the state]
\label{cond:A}
For every $\theta$ and every $t$, the map $X \mapsto f(X, \theta, t)$ is affine:
\[
  f(X_i, \theta, t_i) \;=\; A_i(\theta, t_i)\, X_i \;+\; c_i(\theta, t_i),
\]
with $A_i \in \R^{D\times D}$ and $c_i \in \R^{D}$ \emph{arbitrary} functions of $\theta$ and $t$.
\end{condition}

Note what is not assumed. No linearity, smoothness, boundedness or monotonicity in $\theta$; no
constancy in $t$; and in particular not that $f$ is affine in $(X,\theta)$ jointly, which is a far
stronger requirement and excludes almost every model of interest.

\begin{proposition}
\label{prop:main}
Under Condition~\ref{cond:A}, for each fixed $\theta$ the function $U(\theta,\cdot)$ is an exact
quadratic in $x$, and consequently
\[
  p(x \mid \theta) \;=\; \mathcal N\big(\xstar(\theta),\, \Hxx(\theta)^{-1}\big)
  \quad\text{exactly,}
\]
with
\begin{equation}
\label{eq:Hxx}
  \Hxx(\theta) \;=\; b\,G + W + b\, B(\theta)^\top \Lambda\, B(\theta),
  \qquad B(\theta) = A(\theta) - M,
\end{equation}
\begin{equation}
\label{eq:xstar}
  \xstar(\theta) \;=\; \Hxx(\theta)^{-1}\big[\, b\,G\mu + Wy - b\,B(\theta)^\top \Lambda\, e(\theta)\,\big],
  \qquad e(\theta) = c(\theta) - \dot\mu + M\mu,
\end{equation}
where $A(\theta) = \operatorname{blkdiag}_i(A_i)$ and $c(\theta)$ stacks the $c_i$. Moreover the
elimination of $x$ is exact:
\begin{equation}
\label{eq:exact}
  \int_{\R^{nD}} e^{-U(\theta,x)}\,\mathrm{d}x
  \;=\; e^{-U(\theta,\,\xstar(\theta))}\;(2\pi)^{nD/2} \det \Hxx(\theta)^{-1/2}.
\end{equation}
\end{proposition}

\begin{proof}
Under Condition~\ref{cond:A}, $F(x,\theta) = A(\theta)x + c(\theta)$, so by \eqref{eq:resid}
$r = B(\theta)x + e(\theta)$ is affine in $x$. Then $r^\top\Lambda r$ is quadratic in $x$, and the
remaining $x$-dependent terms of \eqref{eq:U} are quadratic by inspection, so $U(\theta,\cdot)$ is
quadratic. Differentiating,
\[
  \nabla_x U = b\,G(x-\mu) + W(x-y) + b\,B^\top\Lambda\,(Bx + e),
  \qquad
  \nabla^2_x U = b\,G + W + b\,B^\top\Lambda B,
\]
which is \eqref{eq:Hxx} and is independent of $x$; setting $\nabla_x U = 0$ gives
\eqref{eq:xstar}. A quadratic negative log-density is a Gaussian, and \eqref{eq:exact} is the
Gaussian normalising constant, exact rather than asymptotic.
\end{proof}

\subsection{Why $\theta$'s arbitrariness cannot interfere}

$A(\theta)$ and $c(\theta)$ are the only route by which $\theta$ enters, and they enter as
\emph{coefficients}: constants with respect to $x$. They may be arbitrarily wild functions of
$\theta$, and the consequence is that $\Hxx(\theta)$ and $\xstar(\theta)$ are arbitrarily wild
functions of $\theta$. Neither changes the fact that $U(\theta,\cdot)$ is a polynomial of degree
two in $x$.

The general point is that \emph{conditional} Gaussianity places no constraint whatever on how the
conditional's moments depend on the conditioning variable. It is a statement about the shape of one
slice, not about how the slices are arranged.

\begin{remark}[conditional Gaussianity is nearly vacuous as a restriction on the joint]
A model that is Gaussian in $x$ given $\theta$, with arbitrary $\theta$-dependence of the moments,
is a Gaussian \emph{mixture} once $\theta$ is integrated out. Mixtures of Gaussians are dense in
the space of distributions, so Condition~\ref{cond:A} restricts the joint posterior essentially not
at all. What it restricts is the \emph{factorisation}: it says the conditional slices are exactly
Gaussian, which is a computational statement and not a smallness assumption. This is precisely why
the objection in Section~1 feels forceful and is nonetheless answerable.
\end{remark}

\section{Demonstration}
\label{sec:demo}

Take $D = 2$, $p = 3$, and a field affine in $X$ whose $\theta$-dependence is chosen to be as
unpleasant as convenient:
\begin{equation}
\label{eq:demofield}
  f(X, \theta, t) \;=\;
  \begin{pmatrix} -e^{a} & \sin 3b \\ c^{2} & -(1+a^{2})^{-1} \end{pmatrix} X
  \;+\; \begin{pmatrix} \cos 2b \\ e^{-c} \end{pmatrix},
  \qquad \theta = (a,b,c).
\end{equation}
Data were generated at $\theta = (0.3, 0.7, 0.5)$ on a grid of $81$ points over $[0,4]$ with $21$
observations per component and $\sigma = 0.05$, and the mode found by Gauss--Newton in float64.

\paragraph{Is $\Hxx$ independent of $X$?} Proposition~\ref{prop:main} says yes exactly, and
$\|\Hxx(X_1) - \Hxx(X_2)\| / \|\Hxx\|$ at two states sharing one $\theta$ measures the departure.
For \eqref{eq:demofield} it is $0$ to the last bit, at any perturbation size.

\paragraph{Is the quadratic model exact, and how far out?} This is the sharper test, because a
vanishing third derivative at a point is a local statement while \eqref{eq:exact} is global. We
compare $U(\theta, \xstar + \delta)$ against its own second-order expansion
$U(\theta,\xstar) - g^\top\delta + \tfrac12\delta^\top \Hxx \delta$, relative to the change in $U$
itself, along a random direction at growing $\|\delta\|$:

\begin{table}[h]
\centering\small
\begin{tabular}{lrr}
\toprule
$\|\delta\|$ (state units) & \eqref{eq:demofield}, affine in $X$ & FitzHugh--Nagumo, cubic in $X$\\
\midrule
1   & $1.158\times10^{-12}$ & $1.395\times10^{-4}$\\
10  & $1.157\times10^{-12}$ & $3.891\times10^{-2}$\\
100 & $1.157\times10^{-12}$ & $9.479\times10^{-1}$\\
\bottomrule
\end{tabular}
\caption{Relative error of the quadratic model of $U$ in $x$. For the affine field it is flat at
the float64 round-off floor over two orders of magnitude in displacement --- the quadratic is not
an approximation that degrades, it is the function. For FitzHugh--Nagumo it grows until the
quadratic explains essentially nothing.}
\label{tab:taylor}
\end{table}

\paragraph{And in $\theta$?} Along a line in $\theta$ instead, with the same third-difference
stencil, $|U'''|/|U''| = 2.83$ for \eqref{eq:demofield} against $1.07\times10^{-10}$ along a line
in $X$ --- the latter being the round-off floor of a third difference dividing by $h^3 = 10^{-9}$
rather than a measured curvature. The same posterior is exactly quadratic in one block of
coordinates and emphatically not in the other, which is the whole content of this note.

\section{What the condition buys}
\label{sec:buy}

\eqref{eq:exact} is an identity, not an approximation, so the profiled construction of
\cite{companion2},
\begin{equation}
\label{eq:profiled}
  p(\theta) \;\propto\; \pi_\Theta(\theta)\,
  \exp\!\big(-U(\theta, \xstar(\theta))\big)\, \det \Hxx(\theta)^{-1/2},
\end{equation}
is \emph{exact} under Condition~\ref{cond:A}. Three consequences:

\begin{enumerate}[leftmargin=1.6em,itemsep=3pt]
\item \textbf{The dimension reduction is exact.} A $d = p + nD$ dimensional inference problem
becomes a $p$-dimensional one --- $325 \to 3$ on the benchmark of \cite{companion1} --- with no
error incurred in the reduction. Everything that remains is error in the $p$-dimensional integral,
which is a quadrature problem in three to seven dimensions rather than an approximation.
\item \textbf{The mixture representation is exact.} The pipeline reports one Gaussian in $X$ per
$\theta$ node; under Condition~\ref{cond:A} each of those Gaussians is the true conditional, so
every moment computed from the mixture, including $\operatorname{Cov}(\theta, X)$, is exact up to
the quadrature error in $\theta$.
\item \textbf{The inner solve converges in one step.} $\xstar(\theta)$ solves a linear system, so
the inner Gauss--Newton iteration is exact after a single iteration rather than iterating to a
tolerance.
\end{enumerate}

\section{What it does not buy}
\label{sec:notbuy}

The $\theta$-marginal \eqref{eq:profiled} is not Gaussian, is not close to Gaussian, and still has
to be integrated numerically. Both factors carry the arbitrary $\theta$-dependence:
$U(\theta,\xstar(\theta))$ through $A$ and $c$, and $\det\Hxx(\theta)$ through $A$ alone. The
measured $|U'''|/|U''| = 2.83$ of Section~\ref{sec:demo} is that non-Gaussianity, in the same
posterior whose conditional is exactly Gaussian.

Nor does the condition make the mode a good summary. The mode-to-mean displacement that motivates
the whole of \cite{companion1} is a property of the $\theta$-marginal, and Condition~\ref{cond:A}
says nothing about it.

\begin{remark}[a narrower sub-case where more collapses]
If $A$ does not depend on $\theta$ --- $f = AX + c(\theta)$, so the parameters enter only the
forcing --- then $\Hxx$ is constant, the $\det\Hxx(\theta)^{-1/2}$ factor in \eqref{eq:profiled}
becomes an irrelevant constant, and the $\theta$-marginal reduces to
$\pi_\Theta(\theta)\exp(-U(\theta,\xstar(\theta)))$ with $\xstar$ affine in $c(\theta)$. If
additionally $c$ is affine in $\theta$, the whole posterior is jointly Gaussian and there is
nothing left to compute. That is the case the objection of Section~1 was reaching for, and it is
much smaller than Condition~\ref{cond:A}.
\end{remark}

\section{How near are real problems, and a repair to the diagnostic}
\label{sec:near}

$\|\Delta \Hxx\|/\|\Hxx\|$ across two states is a continuous measure, so it reports not only
whether Condition~\ref{cond:A} holds but how far off it is. Measuring that across problems
required fixing it first: the perturbation was an \emph{absolute} displacement, which is not
comparable between a system whose states are $O(1)$ and one whose states reach $10^{5}$. HIV, with
$\|X\|_{\mathrm{rms}} = 1.0\times10^{4}$, was being perturbed by a relative $10^{-5}$ and duly
reported $7.5\times10^{-7}$ --- close enough to zero to suggest it satisfies the condition exactly,
which it does not. The default is now a fixed fraction of the state's own root-mean-square.

\begin{table}[h]
\centering\small
\begin{tabular}{lrrr}
\toprule
system & $\|X\|_{\mathrm{rms}}$ & absolute, $0.7$ & \textbf{relative, $10\%$}\\
\midrule
\eqref{eq:demofield}, affine in $X$ & 0.53 & $0$ & $\mathbf{0}$\\
HIV            & $1.0\times10^{4}$ & $7.5\times10^{-7}$ & $\mathbf{1.1\times10^{-3}}$\\
Lorenz         & 17.1 & $9.1\times10^{-3}$ & $\mathbf{2.2\times10^{-2}}$\\
FitzHugh--Nagumo & 1.16 & $3.8\times10^{-1}$ & $\mathbf{3.2\times10^{-2}}$\\
Hes1           & 1.22 & $4.3\times10^{1}$  & $\mathbf{3.6\times10^{-1}}$\\
\bottomrule
\end{tabular}
\caption{Departure from Condition~\ref{cond:A}. The absolute column is not a like-for-like
comparison and inverts the ordering of HIV against Lorenz and of Lorenz against
FitzHugh--Nagumo. On the relative measure none of the four systems satisfies the condition ---
HIV's mass-action term $\eta T_U V$, Lorenz's $xy$ and $xz$, FitzHugh--Nagumo's cubic and Hes1's
$\exp(X)$ and $(1+P^2)^{-1}$ all violate it --- but they violate it by two and a half orders of
magnitude of difference.}
\label{tab:near}
\end{table}

The ordering in that last column is worth recording, with the caveat that four points support very
little. HIV is nearest the condition, then Lorenz and FitzHugh--Nagumo together, then Hes1 by an
order of magnitude. That is the same ordering produced by every other difficulty measure in this
pipeline: curvature variation over the posterior ($1.21$, $1.9\times10^2$, $3.8\times10^2$,
$1.1\times10^6$), reference-chain divergences ($0\%$, $1.1\%$, $1.4\%$, $13\%$), agreement between
the profiled curvature and the joint Laplace's, and the mode-to-mean distance ($0.16$, $1.42$,
$1.07$, $1.43$ posterior sd). If that is not a coincidence it is unsurprising: departure from
Condition~\ref{cond:A} is exactly the extent to which $\Hxx$ moves as $X$ moves, and a posterior
whose curvature does not move is one every method here finds easy.

\begin{remark}[scope]
Condition~\ref{cond:A} and Proposition~\ref{prop:main} are stated at fixed observation noise. With
$\sigma$ unknown and inferred, \eqref{eq:U} acquires $\sum_d N_d \log \sigma_d^2$ and the $W$ block
becomes $\sigma$-dependent; $p(x \mid \theta, \sigma)$ is still exactly Gaussian at fixed $\sigma$,
but the joint over $(\theta,\sigma)$ is not, and \code{condition\_A} holds $\sigma$ fixed. Worth
knowing before relying on it.
\end{remark}

\section{Summary}

The claim is about $p(X\mid\theta)$ and it is exact. The objection --- that $\theta$ may enter
arbitrarily --- is correct and lands on the $\theta$-marginal, which is not Gaussian and which the
condition says nothing about. What is bought is that eliminating $300$ to $600$ state dimensions
costs nothing in accuracy, leaving a small integral to be done honestly; what is not bought is any
part of that small integral. The models satisfying the condition exactly are linear compartment and
pharmacokinetic systems, and any constant- or time-varying-coefficient linear system --- with rate
constants entering as arbitrarily as one likes, which is the usual case, since they are typically
parameterised on a log scale. None of the four benchmarks here qualifies, though HIV comes within
about $10^{-3}$ of doing so.

\begin{thebibliography}{9}
\small
\bibitem{yang2021} S.~Yang, S.~W.~K.~Wong, S.~C.~Kou.
  Inference of dynamic systems from noisy and sparse data via manifold-constrained Gaussian
  processes. \emph{PNAS} 118(15), 2021.
\bibitem{companion1} A fast, self-certifying Gaussian posterior for MAGI ODE parameter inference.
  Companion technical report, 2026.
\bibitem{companion2} Profiling the states out: posterior inference for MAGI without a Gaussian
  assumption. Companion technical report, 2026.
\end{thebibliography}

\end{document}
